\documentclass{itmo}

\setlist[itemize]{label=---}
\setlist[enumerate]{label=\arabic*)}

\usepackage{dcolumn}
\newcolumntype{d}[1]{D{.}{.}{#1}}

\begin{document}

\makereporttitle{}
	{лабораторной работе №2}
	{Структуры и алгоритмы в базах данных и распределенных системах}
	{Реализация алгоритма геопоиска (kd-tree)}
	{}
	{Постнов~С.~А./P4135}
	{Платонов~А.~В./доцент, Портнов~П.~В./ассистент}

\maketableofcontents

\chapter{Теоретическая часть}

\section{kd--tree}

kd--tree (\textit{k-dimensional tree}) --- это двоичное дерево поиска для точек в $k$-мерном пространстве.
В данной работе используется случай $k=2$ с координатами $(x, y)$.
На каждом уровне дерева выбирается ось разбиения (чередование $x, y, x, y, \ldots$), а в узел помещается точка--медиана по текущей оси.

Базовые операции:
\begin{itemize}
	\item \texttt{insert}: спуск по дереву с чередованием осей и добавление нового листа;
	\item \texttt{find\_nearest}: поиск ближайшей точки к заданной.
\end{itemize}

Для оценки расстояния используется евклидова метрика (квадрат расстояния, без извлечения корня):
$$
d^2((x_1,y_1),(x_2,y_2)) = (x_1-x_2)^2 + (y_1-y_2)^2.
$$

\section{Сложность операций}

В среднем для сбалансированного дерева:
\begin{itemize}
	\item вставка: $O(\log n)$;
	\item поиск ближайшего соседа: около $O(\log n)$, но в худшем случае может деградировать до $O(n)$.
\end{itemize}

В исследовании отдельно измеряются две операции (\texttt{insert} и \texttt{find\_nearest}) и для каждой фиксируются время на операцию и пропускная способность с 95\% доверительным интервалом.

\chapter{Практическая часть}

\section{Реализация структуры}

Публичный интерфейс kd--tree реализован в \texttt{inc/kdtree.h}:
\begin{itemize}
	\item \texttt{new\_tree\_from\_csv} --- построение дерева по CSV с координатами \texttt{x,y};
	\item \texttt{insert} --- добавление узла с точкой;
	\item \texttt{find\_nearest} --- поиск ближайшей точки к заданной $(x,y)$;
	\item \texttt{free\_tree} --- освобождение памяти дерева.
\end{itemize}

Для построения дерева входные точки читаются в массив пар координат, после чего дерево рекурсивно строится по медиане текущего подмассива с чередованием осей разбиения.

\section{Реализация бенчмарка}

Бенчмарк реализован в файле \texttt{src/benchmark.c}.
Генерация входных датасетов выполняется скриптом \texttt{scripts/generate\_bench\_data.py}.

Выходной CSV для каждого датасета содержит:
\begin{itemize}
	\item среднее время вставки на операцию и 95\% ДИ;
	\item среднюю пропускную способность вставки (ops/sec) и 95\% ДИ;
	\item среднее время поиска на операцию и 95\% ДИ;
	\item среднюю пропускную способность поиска (ops/sec) и 95\% ДИ.
\end{itemize}

\chapter{Исследовательская часть}

\section{Аппаратные характеристики}

Все замеры проводились на ноутбуке со следующей конфигурацией:
\begin{itemize}
	\item операционная система: \texttt{macOS 26.3.1} (arm64);
	\item процессор: \texttt{Apple M3 Pro}, 11 ядер;
	\item оперативная память: \texttt{18 ГБ};
	\item компилятор C: \texttt{gcc}, стандарт \texttt{C17}.
\end{itemize}
Во время замеров не выполнялись ресурсоёмкие фоновые задачи, а ноутбук был подключён к сети электропитания для обеспечения максимальной производительности.

\section{Методика замеров}

Для каждого файла из \texttt{data/bench/manifest.txt} выполняются следующие шаги:
\begin{enumerate}
	\item загрузка точек из CSV;
	\item серия замеров \texttt{insert} (вставляются случайные точки в диапазоне координат датасета);
	\item серия замеров \texttt{find\_nearest} (поиск ближайшей точки для координат из датасета).
\end{enumerate}

Параметры замеров:
\begin{itemize}
	\item \texttt{BENCH\_ITERS}=100 итераций (количество итераций);
	\item 100 операций на итерацию для каждой операции.
\end{itemize}

95\% доверительный интервал рассчитывается по формуле:
$$
\Delta_{95} = 1.96 \cdot \sqrt{\frac{s^2}{n}},
$$
где $s^2$ --- выборочная дисперсия по итерациям, $n$ --- число итераций.

\section{Результаты бенчмарка}

В таблице~\ref{tab:kdtree-latency} приведена задержка операций, в таблице~\ref{tab:kdtree-throughput} --- пропускная способность.
\begin{table}[H]
	\centering
	\caption{Задержка операций kd--tree (среднее по итерациям)}
	\label{tab:kdtree-latency}
	\begin{tabularx}{\textwidth}{|l|X|X|}
		\toprule
		$N$ & вставка, нс/оп & поиск, нс/оп \\
		\midrule
		5000 & $122.7 \pm 8.3$ & $100.5 \pm 2.1$ \\
		10000 & $120.8 \pm 2.6$ & $110.3 \pm 3.3$ \\
		20000 & $122.5 \pm 1.8$ & $111.6 \pm 2.5$ \\
		35000 & $128.6 \pm 1.0$ & $128.0 \pm 3.5$ \\
		50000 & $138.3 \pm 2.4$ & $132.7 \pm 2.5$ \\
		75000 & $146.7 \pm 2.1$ & $152.8 \pm 2.4$ \\
		100000 & $165.2 \pm 7.6$ & $154.3 \pm 3.0$ \\
		250000 & $352.0 \pm 17.5$ & $202.7 \pm 7.3$ \\
		500000 & $795.2 \pm 29.2$ & $229.3 \pm 20.9$ \\
		750000 & $1093.2 \pm 43.9$ & $245.6 \pm 29.6$ \\
		1000000 & $1306.7 \pm 53.5$ & $247.9 \pm 31.8$ \\
		\bottomrule
	\end{tabularx}
\end{table}

\begin{table}[H]
	\centering
	\caption{Пропускная способность операций kd--tree (среднее по итерациям)}
	\label{tab:kdtree-throughput}
	\begin{tabularx}{\textwidth}{|l|X|X|}
		\toprule
		$N$ & вставка, млн оп/с & поиск, млн оп/с \\
		\midrule
		5000 & $8.61 \pm 0.28$ & $10.02 \pm 0.14$ \\
		10000 & $8.35 \pm 0.14$ & $9.19 \pm 0.17$ \\
		20000 & $8.20 \pm 0.11$ & $9.03 \pm 0.13$ \\
		35000 & $7.79 \pm 0.06$ & $7.91 \pm 0.14$ \\
		50000 & $7.28 \pm 0.10$ & $7.58 \pm 0.09$ \\
		75000 & $6.84 \pm 0.08$ & $6.57 \pm 0.07$ \\
		100000 & $6.26 \pm 0.18$ & $6.52 \pm 0.09$ \\
		250000 & $2.98 \pm 0.12$ & $5.00 \pm 0.08$ \\
		500000 & $1.30 \pm 0.05$ & $4.54 \pm 0.08$ \\
		750000 & $0.96 \pm 0.04$ & $4.30 \pm 0.08$ \\
		1000000 & $0.80 \pm 0.04$ & $4.28 \pm 0.08$ \\
		\bottomrule
	\end{tabularx}
\end{table}

Из полученных результатов можно сделать следующие выводы:
\begin{itemize}
	\item при малых и средних размерах данных вставка и поиск имеют сопоставимую задержку;
	\item при росте $N$ задержка вставки растёт заметно быстрее, чем задержка поиска;
	\item пропускная способность поиска остаётся более стабильной на больших датасетах.
\end{itemize}

Графики результатов замеров представлены на рисунках ниже.
\includeimage
	{kdtree_bench_avg_latency}
	{f}
	{H}
	{\textwidth}
	{Средняя задержка операций вставки и поиска для kd--tree}

\includeimage
	{kdtree_bench_throughput}
	{f}
	{H}
	{\textwidth}
	{Пропускная способность операций вставки и поиска для kd--tree}

\chapter{Вывод}

В работе реализован kd--tree индекс для геопоиска с поддержкой построения из CSV, динамической вставки и поиска ближайшей точки.
Подготовлен бенчмарк, который измеряет производительность операций \texttt{insert} и \texttt{find\_nearest} на наборах данных различного размера и вычисляет 95\% доверительные интервалы как для времени на операцию, так и для пропускной способности.

Экспериментально показано, что при росте размера данных операция поиска масштабируется стабильнее, чем операция вставки в динамически модифицируемое дерево.
Полученные результаты оформлены в виде таблиц и графиков и могут использоваться для дальнейшей оптимизации реализации (например, улучшение стратегии балансировки и снижение накладных расходов на вставку).

\end{document}
